% TODO: contenido de los anexos A-E
%%%%%%%%%%%%%%%%%%%%%%%%%%%%%%%%%%%%%%%%%%%%%%%%%%%%%%%%%%%%%%%%%%%%%%%%%%%%%%%%%%%%%
%
%  PLANTILLA DE TRABAJO DE GRADO
%  Universidad de San Buenaventura
%  By GHS & PADIA - 2026-2
%
%%%%%%%%%%%%%%%%%%%%%%%%%%%%%%%%%%%%%%%%%%%%%%%%%%%%%%%%%%%%%%%%%%%%%%%%%%%%%%%%%%%%%
%%%%%%%%%%%%%%%%%%%%%%%%%%%%%%%%%%%%%%%%%%%%%%%%%%%%%%%%%%%%%%%%%%%%%%
%  Anexos
%  \anexos (definido en template/sections_config.tex) reinicia el contador
%  y hace que \section numere con letras: Anexo A, Anexo B, ...
%%%%%%%%%%%%%%%%%%%%%%%%%%%%%%%%%%%%%%%%%%%%%%%%%%%%%%%%%%%%%%%%%%%%%%

\clearpage
\phantomsection
\addcontentsline{toc}{section}{ANEXOS}
\section*{ANEXOS}

\anexos


\section{Documento completo SRS (historias de usuario / backlog)}
% Guia: documento completo SRS: historias de usuario y backlog del producto.
\begin{table}[H]
\centering
\caption{Especificación del caso de uso CU-01: Iniciar sesión}
\label{tab:cu01_login}
\renewcommand{\arraystretch}{1.3}
\begin{tabularx}{\textwidth}{|>{\bfseries}p{3.5cm}|X|}
\hline
Elemento & Especificación \\
\hline
ID & CU-01 \\
\hline
Nombre & Iniciar sesión \\
\hline
Actor principal & Usuario registrado \\
\hline
Requisitos relacionados & RF-10 \\
\hline
Objetivo & El objetivo es permitir que los usuarios registrados accedan al sistema utilizando sus credenciales. \\
\hline
Precondiciones & Para poder iniciar sesión, el usuario debe encontrarse previamente registrado en el sistema. \\
\hline
Postcondiciones & El usuario accede exitosamente al sistema y puede acceder a las funcionalidades a las que tenga acceso dependiendo de su rol. \\
\hline
Flujo principal &
\begin{enumerate}
    \item El usuario ingresa su usuario y contraseña.
    \item El sistema valida que la información sea correcta.
    \item El sistema identifica el perfil del usuario.
    \item El sistema permite el acceso a Thary.
\end{enumerate} \\
\hline
Flujos alternativos &
\begin{enumerate}
    \item Si el usuario o la contraseña introducidos son incorrectas, el sistema le informa el usuario que los datos ingresados no son válidos.
    \item Si el usuario no está registrado, el sistema rechaza el acceso y el usuario tiene la oportunidad de registrarse.
\end{enumerate} \\
\hline
\end{tabularx}
\end{table}


\begin{table}[H]
\centering
\caption{Especificación del caso de uso CU-02: Cargar historial de consumo}
\label{tab:cu02_cargar_historial}
\renewcommand{\arraystretch}{1.3}
\begin{tabularx}{\textwidth}{|>{\bfseries}p{3.5cm}|X|}
\hline
Elemento & Especificación \\
\hline
ID & CU-02 \\
\hline
Nombre & Cargar historial de consumo \\
\hline
Actor principal & Administrador \\
\hline
Requisitos relacionados & RF-01, RF-07 \\
\hline
Objetivo & El objetivo del caso de uso es permitirle al Administrador poder cargar un archivo CSV que contenga el historial de consumo de los productos del inventario. \\
\hline
Precondiciones &
\begin{itemize}
    \item Para acceder a esta funcionalidad, el usuario debe haber iniciado sesión con perfil de Administrador.
    \item El archivo debe subirse en formato de CSV.
\end{itemize} \\
\hline
Postcondiciones & El archivo es recibido por el sistema para su posterior procesamiento \\
\hline
Flujo principal &
\begin{enumerate}
    \item El Administrador inicia sesión y es recibido por la pantalla inicial que le da la opción de subir un archivo csv.
    \item El Administrador selecciona un archivo CSV.
    \item El sistema recibe el archivo.
    \item El sistema se encarga de validar el formato y que contenga los campos mínimos requeridos para el funcionamiento del sistema como: fecha, ID del producto y cantidad demandada de cada uno.
    \item El sistema informa que el archivo fue cargado correctamente.
\end{enumerate} \\
\hline
Flujos alternativos &
\begin{enumerate}
    \item Si el archivo no tiene formato CSV, el sistema informa el error con el mensaje "El archivo debe tener extensión .csv"
    \item Si faltan campos obligatorios, el sistema rechaza el archivo e informa cuáles son los datos requeridos, ejemplo: Falta la columna obligatoria: fecha"
    \item Si el usuario no cuenta con los permisos de Administrador, Thary simplemente no le muestra la opción de carga archivo, limitandose con solo mostrarle el dashboard con los resultados y demas pantallas con información.
\end{enumerate} \\
\hline
\end{tabularx}
\end{table}

\begin{table}[H]
\centering
\caption{Especificación del caso de uso CU-03: Procesar archivo}
\label{tab:cu03_procesar_archivo}
\renewcommand{\arraystretch}{1.3}
\begin{tabularx}{\textwidth}{|>{\bfseries}p{3.5cm}|X|}
\hline
Elemento & Especificación \\
\hline
ID & CU-03 \\
\hline
Nombre & Procesar archivo \\
\hline
Actor principal & Administrador \\
\hline
Requisitos relacionados & RF-02, RF-07 \\
\hline
Objetivo & Iniciar el procesamiento del historial de consumo previamente cargado para generar las estimaciones de demanda. \\
\hline
Precondiciones &
\begin{itemize}
    \item El Administrador debe haber iniciado sesión.
    \item Debe existir un archivo de historial cargado y válido.
\end{itemize} \\
\hline
Postcondiciones & El sistema procesa los datos, ejecuta el proceso de predicción y genera los resultados de demanda para los productos. \\
\hline
Flujo principal &
\begin{enumerate}
    \item El Administrador accede al panel de procesamiento.
    \item El sistema verifica que exista un archivo válido.
    \item El Administrador selecciona la opción ``Procesar archivo''.
    \item El sistema procesa y prepara los datos.
    \item El sistema ejecuta el proceso de predicción.
    \item El sistema genera las estimaciones de demanda mensual para los productos.
    \item El sistema genera los resultados necesarios para las etapas posteriores.
\end{enumerate} \\
\hline
Flujos alternativos &
\begin{enumerate}
    \item Si no existe un archivo válido, el sistema informa al Administrador que debe cargar un historial.
    \item Si ocurre un error durante el procesamiento, el sistema informa que no fue posible completar la operación.
    \item Si el usuario no es Administrador, el sistema impide la ejecución.
\end{enumerate} \\
\hline
\end{tabularx}
\end{table}

\begin{table}[H]
\centering
\caption{Especificación del caso de uso CU-04: Consultar demanda pronosticada}
\label{tab:cu04_demanda_pronosticada}
\renewcommand{\arraystretch}{1.3}
\begin{tabularx}{\textwidth}{|>{\bfseries}p{3.5cm}|X|}
\hline
Elemento & Especificación \\
\hline
ID & CU-04 \\
\hline
Nombre & Consultar demanda pronosticada \\
\hline
Actor principal & Personal de Inventario \\
\hline
Requisitos relacionados & RF-02, RF-05 \\
\hline
Objetivo & Permitirle al los empleados del inventario consultar cual es la demanda mensual estimada para cada uno de los producto. \\
\hline
Precondiciones & Como precondición, ya debe de existir un archivo procesado anteriormente por el administrador, también ya deben haber resultados de predicción disponibles. \\
\hline
Postcondiciones & El usuario visualiza las estimaciones de demanda correspondientes a los productos. \\
\hline
Flujo principal &
\begin{enumerate}
    \item El usuario accede al dashboard con las predicciones.
    \item El sistema consulta los resultados generados.
    \item El sistema le muestra al usuario, sin importar si es administrador o usuario, la demanda pronosticada de los productos.
    \item El usuario consulta la información presentada.
\end{enumerate} \\
\hline
Flujos alternativos & Si no existen resultados de predicción, el sistema informa que primero debe realizarse el procesamiento del archivo mediante el siguiente mensaje "Todavía no hay datos cargados. En cuanto tu administrador suba un archivo, vas a poder ver el pronóstico, el inventario y confirmar pedidos." \\
\hline
\end{tabularx}
\end{table}

\begin{table}[H]
\centering
\caption{Especificación del caso de uso CU-05: Consultar estado del inventario}
\label{tab:cu05_estado_inventario}
\renewcommand{\arraystretch}{1.3}
\begin{tabularx}{\textwidth}{|>{\bfseries}p{3.5cm}|X|}
\hline
Elemento & Especificación \\
\hline
ID & CU-05 \\
\hline
Nombre & Consultar estado del inventario \\
\hline
Actor principal & Personal de Inventario \\
\hline
Requisitos relacionados & RF-04, RF-05 \\
\hline
Objetivo & El objetivo es que se pueda consultar el inventario actual de los productos e identificar fácilmente a aquellos que requieran atención o esten en mayor riesgo de desabatecimiento. \\
\hline
Precondiciones & Como precondición para este caso de uso, es importante que ya exista información del inventario disponible y resultados del procesamiento. \\
\hline
Postcondiciones & Las postcondiciones son que el usuario puede visualizar sin problemas el estado actual del inventario y las alertas correspondientes para cada producto. \\
\hline
Flujo principal &
\begin{enumerate}
    \item El usuario accede al módulo de "Estado de inventario".
    \item El sistema consulta la información actual del inventario.
    \item El sistema se encarga de comparar el inventario disponible con el punto de reposición de cada producto.
    \item El sistema identifica los productos que se encuentran en o por debajo de su punto de reposición.
    \item El sistema muestra las alertas visuales correspondientes dependiendo de cada caso, se marcaran de color rojo en caso de estar por debajo del punto de reposición.
\end{enumerate} \\
\hline
Flujos alternativos & Si no existe información de inventario para un producto, el sistema informa que no se dispone de los datos necesarios para determinar su estado. \\
\hline
\end{tabularx}
\end{table}

\begin{table}[H]
\centering
\caption{Especificación del caso de uso CU-06: Consultar recomendaciones de reposición y presupuesto}
\label{tab:cu06_recomendaciones}
\renewcommand{\arraystretch}{1.3}
\begin{tabularx}{\textwidth}{|>{\bfseries}p{3.5cm}|X|}
\hline
Elemento & Especificación \\
\hline
ID & CU-06 \\
\hline
Nombre & Consultar recomendaciones de reposición y presupuesto \\
\hline
Actor principal & Personal de Inventario \\
\hline
Requisitos relacionados & RF-03, RF-05 \\
\hline
Objetivo & Permitir consultar las recomendaciones de reposición generadas para cada producto, teniendo en cuenta distintos datos como la demanda estimada, el inventario que ya se tiene disponible, el tiempo de entrega, que varía de producto a producto, y el presupuesto establecido. \\
\hline
Precondiciones & La principal precondición es que debe existir un archivo csv ya procesado, y el sistema ya debe contar con toda la información necesaria para poder calcular las recomendaciones. \\
\hline
Postcondiciones & El usuario visualiza la cantidad recomendada de compra de una manera clara mediante el uso de gráficas y tablas, también puede ver otros datos relevantes como el punto de reposición y información acerca del presupuesto. \\
\hline
Flujo principal &
\begin{enumerate}
    \item El usuario accede al módulo de recomendaciones, es decir, el dashboard principal.
    \item El sistema consulta las estimaciones de demanda.
    \item El sistema obtiene el inventario actual y los parámetros de reposición.
    \item El sistema considera el tiempo de entrega del proveedor y las restricciones presupuestales.
    \item El sistema calcula la recomendación de reposición.
    \item El sistema le muestra los resultados al usuario.
\end{enumerate} \\
\hline
Flujos alternativos & Si no existen datos suficientes para realizar el cálculo, el sistema informa qué información se encuentra pendiente. \\
\hline
\end{tabularx}
\end{table}


\begin{table}[H]
\centering
\caption{Especificación del caso de uso CU-07: Registrar pedido}
\label{tab:cu07_registrar_pedido}
\renewcommand{\arraystretch}{1.3}
\begin{tabularx}{\textwidth}{|>{\bfseries}p{3.5cm}|X|}
\hline
Elemento & Especificación \\
\hline
ID & CU-07 \\
\hline
Nombre & Registrar pedido \\
\hline
Actor principal & Personal de Inventario \\
\hline
Requisitos relacionados & RF-06 \\
\hline
Objetivo & El objetivo es permitirle al Personal de Inventario registrar un pedido de reposición a partir de las recomendaciones generadas. \\
\hline
Precondiciones & El usuario debe haber iniciado sesión y debe existir información del producto y de la cantidad a solicitar. \\
\hline
Postcondiciones & El pedido queda registrado con la información correspondiente y asociado con el nombre del usuario responsable. \\
\hline
Flujo principal &
\begin{enumerate}
    \item El usuario consulta la página de "Pedidos" del sistema.
    \item Selecciona uno de los pedidos para registrarlos.
    \item Selecciona el proveedor.
    \item A partir de esto, el sistema registra la fecha del pedido.
    \item El sistema identifica al usuario responsable.
    \item El sistema guarda la información del pedido.
    \item El sistema confirma el registro.
\end{enumerate} \\
\hline
Flujos alternativos & Si no se selcciona ningún producto, pero de todas maneras se hace click en el boton para registrar, saldra el siguiente mensaje: No seleccionaste ningún producto para confirmar" \\
\hline
\end{tabularx}
\end{table}

\begin{table}[H]
\centering
\caption{Especificación del caso de uso CU-08: Gestionar usuarios}
\label{tab:cu08_gestionar_usuarios}
\renewcommand{\arraystretch}{1.3}
\begin{tabularx}{\textwidth}{|>{\bfseries}p{3.5cm}|X|}
\hline
Elemento & Especificación \\
\hline
ID & CU-08 \\
\hline
Nombre & Gestionar usuarios \\
\hline
Actor principal & Administrador \\
\hline
Requisitos relacionados & RF-08 \\
\hline
Objetivo & Permitirle al Administrador poder crear cuentas para sus empleados con menos privilegios a los de un administrador. \\
\hline
Precondiciones & El usuario debe haber iniciado sesión con perfil de Administrador. \\
\hline
Postcondiciones & El nuevo empleado queda registrado y puede utilizar el sistema con permisos restringidos. \\
\hline
Flujo principal &
\begin{enumerate}
    \item El Administrador accede a la gestión de usuarios.
    \item Selecciona la opción para crear un usuario.
    \item Ingresa la información requerida.
    \item Define el perfil o nivel de privilegios.
    \item El sistema valida la información.
    \item El sistema registra al nuevo usuario.
\end{enumerate} \\
\hline
Flujos alternativos &
\begin{enumerate}
    \item Si la información es inválida o incompleta, el sistema solicita corregirla mediante un mensaje de error.
    \item Si el usuario ya existe, el sistema informa que no puede crearse una cuenta duplicada.
\end{enumerate} \\
\hline
\end{tabularx}
\end{table}

\begin{table}[H]
\centering
\caption{Especificación del caso de uso CU-09: Deshacer pedido}
\label{tab:cu09_deshacer_pedido}
\renewcommand{\arraystretch}{1.3}
\begin{tabularx}{\textwidth}{|>{\bfseries}p{3.5cm}|X|}
\hline
Elemento & Especificación \\
\hline
ID & CU-09 \\
\hline
Nombre & Deshacer pedido \\
\hline
Actor principal & Personal de Inventario \\
\hline
Requisitos relacionados & RF-10 \\
\hline
Objetivo & El objetivo es permitirle al Personal de Inventario revertir un pedido previamente confirmado, en caso de haberse equivocado o de que la confirmación ya no sea necesaria. \\
\hline
Precondiciones &
\begin{itemize}
    \item El usuario debe haber iniciado sesión.
    \item Debe existir al menos un pedido previamente confirmado por el mismo usuario.
\end{itemize} \\
\hline
Postcondiciones & El pedido queda marcado como no confirmado, y vuelve a aparecer entre las recomendaciones pendientes de confirmar. \\
\hline
Flujo principal &
\begin{enumerate}
    \item El usuario accede a la página de "Pedidos" del sistema.
    \item Consulta el historial de pedidos confirmados.
    \item Selecciona el pedido que desea deshacer.
    \item El sistema solicita confirmación de la acción.
    \item El usuario confirma que desea deshacer el pedido.
    \item El sistema revierte el estado del pedido y lo registra en la bitacora de actividad.
\end{enumerate} \\
\hline
Flujos alternativos &
\begin{enumerate}
    \item Si el usuario cancela la confirmación, el sistema no realiza ningún cambio y el pedido permanece confirmado.
    \item Si el pedido ya fue deshecho previamente por otro usuario, el sistema informa que la acción ya no es posible.
\end{enumerate} \\
\hline
\end{tabularx}
\end{table}


\begin{table}[H]
\centering
\caption{Especificación del caso de uso CU-10: Consultar actividad de la empresa}
\label{tab:cu10_consultar_actividad}
\renewcommand{\arraystretch}{1.3}
\begin{tabularx}{\textwidth}{|>{\bfseries}p{3.5cm}|X|}
\hline
Elemento & Especificación \\
\hline
ID & CU-10 \\
\hline
Nombre & Consultar actividad de la empresa \\
\hline
Actor principal & Administrador \\
\hline
Requisitos relacionados & RF-11 \\
\hline
Objetivo & Permitirle al Administrador consultar la bitacora completa de acciones realizadas por todos los usuarios de su empresa, con fines de auditoria. \\
\hline
Precondiciones &
\begin{itemize}
    \item El usuario debe haber iniciado sesión con perfil de Administrador.
    \item Debe existir al menos un registro de actividad generado previamente en la empresa.
\end{itemize} \\
\hline
Postcondiciones & El administrador visualiza el listado de acciones realizadas por los usuarios de su empresa, ordenadas cronologicamente. \\
\hline
Flujo principal &
\begin{enumerate}
    \item El Administrador accede al módulo de "Actividad".
    \item El sistema consulta el registro de actividad de la empresa del administrador.
    \item El sistema muestra el listado de acciones, indicando el usuario responsable, el tipo de acción y la fecha y hora en que ocurrió.
\end{enumerate} \\
\hline
Flujos alternativos & Si aún no se ha registrado ninguna actividad en la empresa, el sistema informa que no hay eventos para mostrar. \\
\hline
\end{tabularx}
\end{table}


\begin{table}[H]
\centering
\caption{Especificación del caso de uso CU-11: Consultar calendario de pedidos}
\label{tab:cu11_calendario_pedidos}
\renewcommand{\arraystretch}{1.3}
\begin{tabularx}{\textwidth}{|>{\bfseries}p{3.5cm}|X|}
\hline
Elemento & Especificación \\
\hline
ID & CU-11 \\
\hline
Nombre & Consultar calendario de pedidos \\
\hline
Actor principal & Personal de Inventario \\
\hline
Requisitos relacionados & RF-12 \\
\hline
Objetivo & Permitirle al usuario visualizar en un calendario las fechas estimadas en las que debe realizarse el pedido de cada producto, según las recomendaciones generadas por el sistema. \\
\hline
Precondiciones & Debe existir un archivo previamente procesado, con resultados de reposición disponibles. \\
\hline
Postcondiciones & El usuario visualiza, organizadas por fecha, las recomendaciones de pedido proximas a vencer. \\
\hline
Flujo principal &
\begin{enumerate}
    \item El usuario accede al módulo de "Calendario de pedidos".
    \item El sistema consulta las recomendaciones de reposición vigentes.
    \item El sistema organiza dichas recomendaciones según la fecha estimada en la que deberían realizarse.
    \item El sistema muestra el calendario con los pedidos proximos, por producto.
\end{enumerate} \\
\hline
Flujos alternativos & Si no existen recomendaciones de reposición disponibles, el sistema informa que primero debe procesarse un archivo. \\
\hline
\end{tabularx}
\end{table}


\begin{table}[H]
\centering
\caption{Especificación del caso de uso CU-12: Consultar perfil}
\label{tab:cu12_consultar_perfil}
\renewcommand{\arraystretch}{1.3}
\begin{tabularx}{\textwidth}{|>{\bfseries}p{3.5cm}|X|}
\hline
Elemento & Especificación \\
\hline
ID & CU-12 \\
\hline
Nombre & Consultar perfil \\
\hline
Actor principal & Usuario registrado \\
\hline
Requisitos relacionados & RF-13 \\
\hline
Objetivo & Permitirle a cualquier usuario, sin importar su rol, consultar los datos de su propia cuenta y el historial de las acciones que el mismo ha realizado en el sistema. \\
\hline
Precondiciones & El usuario debe haber iniciado sesión. \\
\hline
Postcondiciones & El usuario visualiza la información de su cuenta y sus propios movimientos dentro del sistema. \\
\hline
Flujo principal &
\begin{enumerate}
    \item El usuario accede al módulo de "Perfil".
    \item El sistema consulta los datos de la cuenta del usuario.
    \item El sistema consulta los movimientos propios del usuario (pedidos confirmados o revertidos, y si es administrador, también sus archivos cargados y los usuarios que ha creado).
    \item El sistema muestra la información recopilada.
\end{enumerate} \\
\hline
Flujos alternativos & Si el usuario aún no ha realizado ninguna acción registrada en el sistema, el sistema muestra los datos de la cuenta sin movimientos asociados. \\
\hline
\end{tabularx}
\end{table}

\section{Modelos de datos extensos, diccionarios de datos, mockups}
% Guía: modelos de datos extensos, diccionarios de datos y mockups.

Este anexo presenta el diccionario de datos completo del esquema relacional de Thary, complementando la descripción resumida presentada en el capítulo 3. Todas las tablas se implementan sobre MySQL, con codificación \texttt{utf8mb4} y motor de almacenamiento InnoDB.

\subsection*{Tabla \texttt{empresas}}
\begin{table}[H]
\centering
\renewcommand{\arraystretch}{1.3}
\begin{tabularx}{\textwidth}{|>{\raggedright\arraybackslash}p{3.0cm}|l|>{\raggedright\arraybackslash}p{3.2cm}|X|}
\hline
\textbf{Campo} & \textbf{Tipo} & \textbf{Restricción} & \textbf{Descripción} \\
\hline
id & VARCHAR(32) & PRIMARY KEY & Identificador único de la organización \\
\hline
nombre\_\discretionary{}{}{}empresa & VARCHAR(120) & & Nombre de la organización, visible en el encabezado \\
\hline
fecha\_\discretionary{}{}{}creación & DATETIME & DEFAULT CURRENT\_\discretionary{}{}{}TIMESTAMP & Fecha de registro de la empresa \\
\hline
\end{tabularx}
\end{table}

\subsection*{Tabla \texttt{usuarios}}
\begin{table}[H]
\centering
\renewcommand{\arraystretch}{1.3}
\begin{tabularx}{\textwidth}{|>{\raggedright\arraybackslash}p{3.0cm}|l|>{\raggedright\arraybackslash}p{3.2cm}|X|}
\hline
\textbf{Campo} & \textbf{Tipo} & \textbf{Restricción} & \textbf{Descripción} \\
\hline
id & INT & PRIMARY KEY, AUTO\_\discretionary{}{}{}INCREMENT & Identificador único del usuario \\
\hline
nombre & VARCHAR(120) & NOT NULL & Nombre completo del usuario \\
\hline
username & VARCHAR(80) & NOT NULL, UNIQUE & Nombre de usuario para inicio de sesión \\
\hline
email & VARCHAR(160) & NOT NULL, UNIQUE & Correo electrónico del usuario \\
\hline
password\_\discretionary{}{}{}hash & VARCHAR(255) & NOT NULL & Hash de la contraseña, generado con \textit{scrypt} \\
\hline
rol & ENUM('admin','empleado') & NOT NULL, DEFAULT 'empleado' & Rol del usuario dentro de la empresa \\
\hline
empresa\_\discretionary{}{}{}id & VARCHAR(32) & NOT NULL, FK $\to$ empresas.id & Empresa a la que pertenece el usuario \\
\hline
\end{tabularx}
\end{table}

\subsection*{Tabla \texttt{pedidos}}
\begin{table}[H]
\centering
\renewcommand{\arraystretch}{1.3}
\begin{tabularx}{\textwidth}{|>{\raggedright\arraybackslash}p{3.0cm}|l|>{\raggedright\arraybackslash}p{3.2cm}|X|}
\hline
\textbf{Campo} & \textbf{Tipo} & \textbf{Restricción} & \textbf{Descripción} \\
\hline
id & INT & PRIMARY KEY, AUTO\_\discretionary{}{}{}INCREMENT & Identificador único del pedido \\
\hline
empresa\_\discretionary{}{}{}id & VARCHAR(32) & NOT NULL, FK $\to$ empresas.id & Empresa a la que pertenece el pedido \\
\hline
producto\_\discretionary{}{}{}id & VARCHAR(64) & NOT NULL & Identificador del producto pedido \\
\hline
proveedor & VARCHAR(160) & & Proveedor elegido para el pedido \\
\hline
cantidad & DOUBLE & & Cantidad confirmada \\
\hline
costo & DOUBLE & & Costo estimado del pedido \\
\hline
usuario & VARCHAR(120) & & Usuario responsable de confirmar el pedido \\
\hline
fecha\_\discretionary{}{}{}hora & DATETIME & & Fecha y hora de confirmación \\
\hline
\multicolumn{4}{l}{\textit{Clave única: (empresa\_id, producto\_id)}} \\
\hline
\end{tabularx}
\end{table}

\subsection*{Tabla \texttt{actividad}}
\begin{table}[H]
\centering
\renewcommand{\arraystretch}{1.3}
\begin{tabularx}{\textwidth}{|>{\raggedright\arraybackslash}p{3.0cm}|l|>{\raggedright\arraybackslash}p{3.2cm}|X|}
\hline
\textbf{Campo} & \textbf{Tipo} & \textbf{Restricción} & \textbf{Descripción} \\
\hline
id & BIGINT & PRIMARY KEY, AUTO\_\discretionary{}{}{}INCREMENT & Identificador único del evento \\
\hline
empresa\_\discretionary{}{}{}id & VARCHAR(32) & NOT NULL, FK $\to$ empresas.id & Empresa asociada al evento \\
\hline
username & VARCHAR(80) & NOT NULL & Usuario que realizó la acción \\
\hline
nombre\_\discretionary{}{}{}usuario & VARCHAR(120) & NOT NULL & Nombre del usuario que realizó la acción \\
\hline
tipo & VARCHAR(30) & NOT NULL & Tipo de acción (archivo\_subido, pedido\_confirmado, etc.) \\
\hline
descripcion & VARCHAR(300) & NOT NULL & Descripción legible del evento \\
\hline
fecha & DATETIME & DEFAULT CURRENT\_\discretionary{}{}{}TIMESTAMP & Fecha y hora del evento \\
\hline
\end{tabularx}
\end{table}

\subsection*{Tabla \texttt{corridas}}
\begin{table}[H]
\centering
\renewcommand{\arraystretch}{1.3}
\begin{tabularx}{\textwidth}{|>{\raggedright\arraybackslash}p{3.0cm}|l|>{\raggedright\arraybackslash}p{3.2cm}|X|}
\hline
\textbf{Campo} & \textbf{Tipo} & \textbf{Restricción} & \textbf{Descripción} \\
\hline
id & INT & PRIMARY KEY, AUTO\_\discretionary{}{}{}INCREMENT & Identificador de la corrida (tabla padre) \\
\hline
empresa\_\discretionary{}{}{}id & VARCHAR(32) & NOT NULL, FK $\to$ empresas.id & Empresa que ejecutó la corrida \\
\hline
fecha\_\discretionary{}{}{}corrida & DATETIME & DEFAULT CURRENT\_\discretionary{}{}{}TIMESTAMP & Fecha de la corrida \\
\hline
horizonte\_\discretionary{}{}{}meses & INT & & Meses hacia adelante pronosticados \\
\hline
presupuesto\_\discretionary{}{}{}capital\_\discretionary{}{}{}trabajo & DOUBLE & & Presupuesto total disponible \\
\hline
presupuesto\_\discretionary{}{}{}por\_\discretionary{}{}{}producto & DOUBLE & & Presupuesto asignado por producto \\
\hline
mape\_\discretionary{}{}{}final & DOUBLE & & MAPE del sistema final \\
\hline
wape\_\discretionary{}{}{}final & DOUBLE & & WAPE del sistema final \\
\hline
costo\_\discretionary{}{}{}total\_\discretionary{}{}{}pedidos\_\discretionary{}{}{}urgentes & DOUBLE & & Costo total de productos en alerta \\
\hline
valor\_\discretionary{}{}{}total\_\discretionary{}{}{}inventario & DOUBLE & & Valor total del inventario en la corrida \\
\hline
productos\_\discretionary{}{}{}sin\_\discretionary{}{}{}dato\_\discretionary{}{}{}inventario & INT & & Cantidad de productos sin inventario registrado \\
\hline
productos\_\discretionary{}{}{}en\_\discretionary{}{}{}alerta & INT & & Cantidad de productos en o bajo el punto de reorden \\
\hline
features\_\discretionary{}{}{}regresión\_\discretionary{}{}{}lineal & TEXT & & Variables seleccionadas por RFE/VIF (serializado) \\
\hline
\end{tabularx}
\end{table}

\subsection*{Tabla \texttt{segmentacion}}
\begin{table}[H]
\centering
\renewcommand{\arraystretch}{1.3}
\begin{tabularx}{\textwidth}{|>{\raggedright\arraybackslash}p{3.0cm}|l|>{\raggedright\arraybackslash}p{3.2cm}|X|}
\hline
\textbf{Campo} & \textbf{Tipo} & \textbf{Restricción} & \textbf{Descripción} \\
\hline
id & BIGINT & PRIMARY KEY, AUTO\_\discretionary{}{}{}INCREMENT & Identificador único \\
\hline
corrida\_\discretionary{}{}{}id & INT & NOT NULL, FK $\to$ corridas.id (CASCADE) & Corrida a la que pertenece \\
\hline
producto\_\discretionary{}{}{}id & VARCHAR(64) & & Identificador del producto \\
\hline
nombre\_\discretionary{}{}{}producto & VARCHAR(200) & & Nombre del producto \\
\hline
categoría & VARCHAR(120) & & Categoría o rubro del producto \\
\hline
demanda\_\discretionary{}{}{}total & DOUBLE & & Demanda total histórica \\
\hline
media & DOUBLE & & Demanda promedio mensual \\
\hline
desviación & DOUBLE & & Desviación estándar de la demanda \\
\hline
cv & DOUBLE & & Coeficiente de variación \\
\hline
variabilidad & VARCHAR(20) & & Clasificación Alta/Media/Baja \\
\hline
demanda\_\discretionary{}{}{}acumulada\_\discretionary{}{}{}pct & DOUBLE & & Porcentaje acumulado de demanda \\
\hline
categoría\_\discretionary{}{}{}abc & VARCHAR(5) & & Clasificación A/B/C \\
\hline
método\_\discretionary{}{}{}campeón & VARCHAR(40) & & Modelo seleccionado como campeón \\
\hline
mae\_\discretionary{}{}{}media\_\discretionary{}{}{}móvil & DOUBLE & & MAE de backtest del modelo Media Móvil para este producto \\
\hline
mae\_\discretionary{}{}{}regresión\_\discretionary{}{}{}lineal & DOUBLE & & MAE de backtest del modelo Regresión Lineal para este producto \\
\hline
mae\_\discretionary{}{}{}xgboost & DOUBLE & & MAE de backtest del modelo XGBoost para este producto \\
\hline
proveedor\_\discretionary{}{}{}principal & VARCHAR(160) & & Proveedor principal del producto \\
\hline
proveedor\_\discretionary{}{}{}alterno & VARCHAR(160) & & Proveedor alterno del producto \\
\hline
\end{tabularx}
\end{table}


\subsection*{Tabla \texttt{metricas\_modelos}}
\begin{table}[H]
\centering
\renewcommand{\arraystretch}{1.3}
\begin{tabularx}{\textwidth}{|>{\raggedright\arraybackslash}p{3.0cm}|l|>{\raggedright\arraybackslash}p{3.2cm}|X|}
\hline
\textbf{Campo} & \textbf{Tipo} & \textbf{Restricción} & \textbf{Descripción} \\
\hline
id & BIGINT & PRIMARY KEY, AUTO\_\discretionary{}{}{}INCREMENT & Identificador único \\
\hline
corrida\_\discretionary{}{}{}id & INT & NOT NULL, FK $\to$ corridas.id (CASCADE) & Corrida a la que pertenece \\
\hline
modelo & VARCHAR(60) & & Nombre del modelo evaluado \\
\hline
mae & DOUBLE & & Error absoluto medio \\
\hline
rmse & DOUBLE & & Raíz del error cuadrático medio \\
\hline
r2 & DOUBLE & & Coeficiente de determinación \\
\hline
mediana\_\discretionary{}{}{}error & DOUBLE & & Mediana del error absoluto \\
\hline
mape & DOUBLE & & Error porcentual absoluto medio \\
\hline
wape & DOUBLE & & Error porcentual absoluto ponderado \\
\hline
\end{tabularx}
\end{table}

\subsection*{Tabla \texttt{predicciones\_backtest}}
\begin{table}[H]
\centering
\renewcommand{\arraystretch}{1.3}
\begin{tabularx}{\textwidth}{|>{\raggedright\arraybackslash}p{3.0cm}|l|>{\raggedright\arraybackslash}p{3.2cm}|X|}
\hline
\textbf{Campo} & \textbf{Tipo} & \textbf{Restricción} & \textbf{Descripción} \\
\hline
id & BIGINT & PRIMARY KEY, AUTO\_\discretionary{}{}{}INCREMENT & Identificador único \\
\hline
corrida\_\discretionary{}{}{}id & INT & NOT NULL, FK $\to$ corridas.id (CASCADE) & Corrida a la que pertenece \\
\hline
fecha & VARCHAR(10) & & Fecha del mes evaluado (AAAA-MM-DD) \\
\hline
producto\_\discretionary{}{}{}id & VARCHAR(64) & & Identificador del producto \\
\hline
demanda & DOUBLE & & Demanda real observada \\
\hline
categoría\_\discretionary{}{}{}abc & VARCHAR(5) & & Clasificación ABC del producto \\
\hline
variabilidad & VARCHAR(20) & & Clasificación de variabilidad \\
\hline
predicción\_\discretionary{}{}{}media\_\discretionary{}{}{}móvil & DOUBLE & & Predicción del modelo Media Móvil \\
\hline
predicción\_\discretionary{}{}{}regresión\_\discretionary{}{}{}lineal & DOUBLE & & Predicción del modelo Regresión Lineal \\
\hline
predicción\_\discretionary{}{}{}xgboost & DOUBLE & & Predicción del modelo XGBoost \\
\hline
predicción\_\discretionary{}{}{}final & DOUBLE & & Predicción del método campeón seleccionado \\
\hline
error\_\discretionary{}{}{}final & DOUBLE & & Error de la predicción final \\
\hline
método\_\discretionary{}{}{}usado & VARCHAR(40) & & Modelo seleccionado para ese producto y mes \\
\hline
\end{tabularx}
\end{table}

\subsection*{Tabla \texttt{pronostico\_futuro}}
\begin{table}[H]
\centering
\renewcommand{\arraystretch}{1.3}
\begin{tabularx}{\textwidth}{|>{\raggedright\arraybackslash}p{3.0cm}|l|>{\raggedright\arraybackslash}p{3.2cm}|X|}
\hline
\textbf{Campo} & \textbf{Tipo} & \textbf{Restricción} & \textbf{Descripción} \\
\hline
id & BIGINT & PRIMARY KEY, AUTO\_\discretionary{}{}{}INCREMENT & Identificador único \\
\hline
corrida\_\discretionary{}{}{}id & INT & NOT NULL, FK $\to$ corridas.id (CASCADE) & Corrida a la que pertenece \\
\hline
producto\_\discretionary{}{}{}id & VARCHAR(64) & & Identificador del producto \\
\hline
nombre\_\discretionary{}{}{}producto & VARCHAR(200) & & Nombre del producto \\
\hline
categoría & VARCHAR(120) & & Categoría del producto \\
\hline
mes\_\discretionary{}{}{}pronosticado & VARCHAR(7) & & Mes futuro pronosticado (AAAA-MM) \\
\hline
horizonte & INT & & Número de mes dentro del horizonte \\
\hline
predicción\_\discretionary{}{}{}xgboost & DOUBLE & & Predicción de XGBoost \\
\hline
predicción\_\discretionary{}{}{}regresión\_\discretionary{}{}{}lineal & DOUBLE & & Predicción de Regresión Lineal \\
\hline
predicción\_\discretionary{}{}{}baseline & DOUBLE & & Predicción de Media Móvil \\
\hline
predicción\_\discretionary{}{}{}final & DOUBLE & & Predicción del método campeón \\
\hline
predicción\_\discretionary{}{}{}min & DOUBLE & & Límite inferior del rango de la predicción final (predicción menos el MAE de backtest del método usado) \\
\hline
predicción\_\discretionary{}{}{}max & DOUBLE & & Límite superior del rango de la predicción final \\
\hline
método\_\discretionary{}{}{}usado & VARCHAR(40) & & Modelo seleccionado para ese producto \\
\hline
\end{tabularx}
\end{table}


\subsection*{Tabla \texttt{reorder}}
\begin{table}[H]
\centering
\renewcommand{\arraystretch}{1.3}
\begin{tabularx}{\textwidth}{|>{\raggedright\arraybackslash}p{3.0cm}|l|>{\raggedright\arraybackslash}p{3.2cm}|X|}
\hline
\textbf{Campo} & \textbf{Tipo} & \textbf{Restricción} & \textbf{Descripción} \\
\hline
id & BIGINT & PRIMARY KEY, AUTO\_\discretionary{}{}{}INCREMENT & Identificador único \\
\hline
corrida\_\discretionary{}{}{}id & INT & NOT NULL, FK $\to$ corridas.id (CASCADE) & Corrida a la que pertenece \\
\hline
producto\_\discretionary{}{}{}id & VARCHAR(64) & & Identificador del producto \\
\hline
nombre\_\discretionary{}{}{}producto & VARCHAR(200) & & Nombre del producto \\
\hline
categoría & VARCHAR(120) & & Categoría del producto \\
\hline
proveedor\_\discretionary{}{}{}principal & VARCHAR(160) & & Proveedor principal \\
\hline
proveedor\_\discretionary{}{}{}alterno & VARCHAR(160) & & Proveedor alterno \\
\hline
categoría\_\discretionary{}{}{}abc & VARCHAR(5) & & Clasificación ABC \\
\hline
variabilidad & VARCHAR(20) & & Clasificación de variabilidad \\
\hline
método\_\discretionary{}{}{}pronóstico & VARCHAR(40) & & Modelo usado para pronosticar \\
\hline
demanda\_\discretionary{}{}{}promedio\_\discretionary{}{}{}mensual\_\discretionary{}{}{}pronosticada & DOUBLE & & Demanda promedio pronosticada \\
\hline
lead\_\discretionary{}{}{}time\_\discretionary{}{}{}semanas & DOUBLE & & Tiempo de entrega en semanas \\
\hline
lead\_\discretionary{}{}{}time\_\discretionary{}{}{}meses & DOUBLE & & Tiempo de entrega en meses \\
\hline
stock\_\discretionary{}{}{}seguridad & DOUBLE & & Stock de seguridad calculado (supuesto de demanda normal) \\
\hline
punto\_\discretionary{}{}{}reorden & DOUBLE & & Punto de reorden calculado (supuesto de demanda normal) \\
\hline
stock\_\discretionary{}{}{}seguridad\_\discretionary{}{}{}lognormal & DOUBLE & & Stock de seguridad calculado asumiendo demanda log-normal, a modo de comparación \\
\hline
punto\_\discretionary{}{}{}reorden\_\discretionary{}{}{}lognormal & DOUBLE & & Punto de reorden calculado asumiendo demanda log-normal, a modo de comparación \\
\hline
inventario\_\discretionary{}{}{}actual & DOUBLE & & Inventario disponible actual \\
\hline
eoq & DOUBLE & & Cantidad económica de pedido \\
\hline
límite\_\discretionary{}{}{}presupuesto\_\discretionary{}{}{}unidades & DOUBLE & & Límite de unidades según presupuesto \\
\hline
cantidad\_\discretionary{}{}{}sugerida\_\discretionary{}{}{}pedido & DOUBLE & & Cantidad final sugerida \\
\hline
costo\_\discretionary{}{}{}unitario & DOUBLE & & Costo unitario del producto \\
\hline
costo\_\discretionary{}{}{}estimado\_\discretionary{}{}{}pedido & DOUBLE & & Costo total estimado del pedido \\
\hline
valor\_\discretionary{}{}{}inventario\_\discretionary{}{}{}actual & DOUBLE & & Valor del inventario disponible \\
\hline
margen\_\discretionary{}{}{}unitario & DOUBLE & & Margen unitario del producto (precio menos costo unitario) \\
\hline
costo\_\discretionary{}{}{}ruptura\_\discretionary{}{}{}estimado & DOUBLE & & Costo estimado de ruptura (\textit{stockout}) \\
\hline
fecha\_\discretionary{}{}{}estimada\_\discretionary{}{}{}pedido & VARCHAR(10) & & Fecha estimada en que debería pedirse (AAAA-MM-DD) \\
\hline
dias\_\discretionary{}{}{}para\_\discretionary{}{}{}pedido & INT & & Días restantes hasta la fecha estimada \\
\hline
ordenar & VARCHAR(5) & & Indicador (``SI''/``NO'') de alerta de reposición \\
\hline
\end{tabularx}
\end{table}

\subsection*{Tabla \texttt{validacion\_cruzada\_folds}}
\begin{table}[H]
\centering
\renewcommand{\arraystretch}{1.3}
\begin{tabularx}{\textwidth}{|>{\raggedright\arraybackslash}p{3.0cm}|l|>{\raggedright\arraybackslash}p{3.2cm}|X|}
\hline
\textbf{Campo} & \textbf{Tipo} & \textbf{Restricción} & \textbf{Descripción} \\
\hline
id & BIGINT & PRIMARY KEY, AUTO\_\discretionary{}{}{}INCREMENT & Identificador único \\
\hline
corrida\_\discretionary{}{}{}id & INT & NOT NULL, FK $\to$ corridas.id (CASCADE) & Corrida a la que pertenece \\
\hline
fold & INT & & Número de fold (1, 2, 3...) \\
\hline
meses\_\discretionary{}{}{}entrenamiento & INT & & Meses usados para entrenar en ese fold \\
\hline
periodo\_\discretionary{}{}{}test & VARCHAR(30) & & Periodo evaluado en ese fold \\
\hline
modelo & VARCHAR(40) & & Modelo evaluado (incluye ``Sistema final'') \\
\hline
mae & DOUBLE & & MAE obtenido en ese fold \\
\hline
rmse & DOUBLE & & RMSE obtenido en ese fold \\
\hline
mape & DOUBLE & & MAPE obtenido en ese fold \\
\hline
wape & DOUBLE & & WAPE obtenido en ese fold \\
\hline
\end{tabularx}
\end{table}

\subsection*{Tabla \texttt{validacion\_cruzada\_resumen}}
\begin{table}[H]
\centering
\renewcommand{\arraystretch}{1.3}
\begin{tabularx}{\textwidth}{|>{\raggedright\arraybackslash}p{3.0cm}|l|>{\raggedright\arraybackslash}p{3.2cm}|X|}
\hline
\textbf{Campo} & \textbf{Tipo} & \textbf{Restricción} & \textbf{Descripción} \\
\hline
id & BIGINT & PRIMARY KEY, AUTO\_\discretionary{}{}{}INCREMENT & Identificador único \\
\hline
corrida\_\discretionary{}{}{}id & INT & NOT NULL, FK $\to$ corridas.id (CASCADE) & Corrida a la que pertenece \\
\hline
modelo & VARCHAR(40) & & Modelo resumido \\
\hline
métrica & VARCHAR(20) & & Métrica resumida (MAE, RMSE, MAPE, WAPE) \\
\hline
promedio & DOUBLE & & Promedio de la métrica entre folds \\
\hline
desvio & DOUBLE & & Desviación estándar entre folds \\
\hline
folds & INT & & Número de folds considerados \\
\hline
\end{tabularx}
\end{table}

\section{Manual de usuario y/o guía de instalación y despliegue}
% Guía: manual de usuario y guía de instalación y despliegue.

\subsection*{Instalación en entorno local}

\begin{enumerate}
    \item Clonar el repositorio del proyecto y acceder a la carpeta creada:
    \begin{minted}{bash}
git clone https://github.com/danna0723/Thary_Proyecto.git
cd Thary_Proyecto
    \end{minted}
    Se requiere tener instalado Python 3.12 o superior.
    \item Crear un entorno virtual e instalar las dependencias:
    \begin{minted}{bash}
python -m venv venv
venv\Scripts\activate       # En Windows
pip install -r requirements.txt
    \end{minted}
    \item Tener un servidor MySQL corriendo localmente (o accesible mediante las variables del siguiente paso) y un usuario con permisos sobre la base de datos del sistema. Por ejemplo:
    \begin{minted}{sql}
CREATE USER 'thary_app'@'localhost' IDENTIFIED BY '<su_contraseña>';
GRANT ALL PRIVILEGES ON sistema_inventario.* TO 'thary_app'@'localhost';
    \end{minted}
    Al iniciar, el sistema crea automáticamente la base de datos (si el usuario tiene permiso para hacerlo) y todas las tablas necesarias.
    \item Configurar las variables de entorno para la conexión a MySQL (\texttt{DB\_HOST}, \texttt{DB\_PORT}, \texttt{DB\_USER}, \texttt{DB\_PASSWORD}, \texttt{DB\_NAME}). Si no se definen, el sistema usa los valores \texttt{127.0.0.1}, puerto \texttt{3307}, usuario \texttt{thary\_app} y base de datos \texttt{sistema\_inventario}, con contraseña vacía. Por ejemplo, en PowerShell:
    \begin{minted}{powershell}
$env:DB_PASSWORD = "<su_contraseña>"
    \end{minted}
    Las variables definidas de esta forma solo duran mientras la terminal esté abierta.
    \item Ejecutar la aplicación. En este modo de desarrollo no hace falta definir \texttt{SECRET\_KEY}, porque se usa una clave de desarrollo por defecto:
    \begin{minted}{bash}
python run.py
    \end{minted}
    \item Acceder desde el navegador a \texttt{http://localhost:5000}.
\end{enumerate}

\subsection*{Despliegue en producción (Railway)}

\begin{enumerate}
    \item Crear un proyecto en Railway y agregar dos servicios: la aplicación web (conectada al repositorio de GitHub) y una instancia de MySQL.
    \item Configurar las variables de entorno del servicio web:
    \begin{itemize}
        \item Las de conexión a la base de datos (\texttt{DB\_HOST}, \texttt{DB\_PORT}, \texttt{DB\_USER}, \texttt{DB\_PASSWORD}, \texttt{DB\_NAME}), apuntando a las credenciales del servicio MySQL de Railway.
        \item \texttt{SECRET\_KEY}: una cadena aleatoria larga que se usa para firmar las sesiones de los usuarios. Sin esta variable la aplicación no arranca, y no debe compartirse ni escribirse en el código.
    \end{itemize}
    Los nombres de las variables no deben tener espacios al inicio ni al final, porque eso hace fallar la construcción del servicio.
    \item El archivo \texttt{Procfile} ya define el comando de arranque en producción mediante Gunicorn:
    \begin{minted}{text}
web: gunicorn run:app --bind 0.0.0.0:$PORT --workers 2 --timeout 120
    \end{minted}
    \item Railway detecta automáticamente la versión de Python especificada en \texttt{.python-version} e instala las dependencias de \texttt{requirements.txt}.
    \item Al iniciar, la aplicación crea automáticamente las tablas necesarias en la base de datos (función \texttt{crear\_tablas()}).
    \item Se recomienda activar un volumen persistente para el servicio web y la funcionalidad de \textit{backups} automáticos para el servicio de MySQL (Sección 4.7.1).
\end{enumerate}

\subsection*{Manual de usuario}

\subsubsection*{Administrador}
\begin{enumerate}
    \item Registrarse: quien crea una cuenta desde la pantalla de registro obtiene una empresa propia y queda como su administrador. Después, iniciar sesión.
    \item Desde el menú lateral, entrar a ``Subir archivo''. Allí se indica el número de meses a pronosticar y el presupuesto disponible, y se carga un archivo CSV con el historial de consumo. Las columnas obligatorias son la fecha, el código del producto y la cantidad vendida; las opcionales habilitan funciones adicionales (por ejemplo, el inventario activa las alertas de pedido). La misma pantalla incluye el detalle de las columnas que reconoce el sistema y una plantilla de ejemplo descargable.
    \item Seleccionar ``Procesar archivo'' para ejecutar el pipeline de predicción.
    \item Consultar el ``Panel técnico'' para revisar la comparación de modelos y la segmentación de productos.
    \item Crear cuentas nuevas para el personal de inventario desde ``Usuarios''.
    \item Consultar la bitácora de actividad de la empresa desde ``Actividad''.
\end{enumerate}

\subsubsection*{Personal de inventario}
\begin{enumerate}
    \item Iniciar sesión con las credenciales proporcionadas por el administrador.
    \item Consultar, dentro de ``Análisis'', el pronóstico de demanda, el estado del inventario, el presupuesto y costos, y el calendario de pedidos. El ícono de la campana, en la parte superior, indica cuántos productos necesitan pedido.
    \item Descargar el pronóstico de demanda en formato CSV desde la pantalla de pronóstico, con la opción de elegir un rango de meses.
    \item Confirmar un pedido sugerido desde la pantalla ``Pedido'', seleccionando el proveedor correspondiente.
    \item Deshacer un pedido confirmado por error con el botón ``Deshacer'' de la misma pantalla.
    \item Consultar su propio perfil y el historial de sus acciones desde ``Mi perfil''.
\end{enumerate}


\section{Enlaces a repositorios de código fuente y tableros de gestión (Azure DevOps, Jira, GitHub)}
% Guía: enlaces a repositorios de código fuente y tableros de gestión.

El código fuente de Thary se gestionó con el sistema de control de versiones Git, y el repositorio se alojó en GitHub. A continuación se presentan los enlaces a los recursos del proyecto.

\begin{table}[H]
\centering
\renewcommand{\arraystretch}{1.3}
\begin{tabularx}{\textwidth}{|l|X|}
\hline
\textbf{Recurso} & \textbf{Enlace} \\
\hline
Repositorio de código fuente (GitHub) & \url{https://github.com/danna0723/Thary_Proyecto} \\
\hline
Flujo de integración continua (GitHub Actions) & \url{https://github.com/danna0723/Thary_Proyecto/actions} \\
\hline
Historial de cambios (commits) & \url{https://github.com/danna0723/Thary_Proyecto/commits/main} \\
\hline
Sistema desplegado (Railway) & \url{https://web-production-6ac7e5.up.railway.app} \\
\hline
\end{tabularx}
\end{table}

\subsection*{Repositorio de código fuente}

El repositorio contiene la aplicación completa (código en \texttt{app/}), las pruebas automatizadas (\texttt{tests/}), la definición del flujo de integración continua (\texttt{.github/workflows/tests.yml}) y los archivos de configuración del despliegue (\texttt{Procfile} y \texttt{requirements.txt}). El desarrollo se realizó sobre una única rama principal (\texttt{main}), con un historial de 16 commits entre el 14 y el 26 de septiembre de 2026 al momento de redactar este documento.

\subsection*{Integración y despliegue continuos}

Cada envío de cambios (\textit{push}) o \textit{pull request} al repositorio ejecuta automáticamente las pruebas automatizadas mediante GitHub Actions, con Python 3.12. Además, el servicio en Railway está conectado al repositorio de GitHub, de modo que cada cambio enviado a la rama principal inicia un nuevo despliegue.

\subsection*{Tableros de gestión}

% Elige UNA de estas dos opciones y borra la otra:

% OPCIÓN A (si NO usaste un tablero externo):
Para el seguimiento del desarrollo del sistema no se utilizó alguna herramienta externa de gestión de proyectos como Jira o Azure DevOps, ya que el avance del desarrollo se registró mediante el historial de commits en GitHub, cuyos mensajes describen los cambios realizados en cada etapa.

% OPCIÓN B (si SÍ usaste un tablero):
% El seguimiento de las tareas del proyecto se realizó en \textit{<nombre de la herramienta>}, disponible en \url{<enlace al tablero>}.


\section{Protocolo y seguimiento de la revisión sistemática de literatura (si aplica)}
% Guia: protocolo y seguimiento de la revision sistematica de literatura.
Este anexo no aplica al presente trabajo. La revisión de literatura presentada en el Capítulo 2 se realizo mediante un proceso de busqueda y selección tradicional, con los criterios de inclusión y exclusion descritos en dicho caítulo, pero sin seguir un protocolo formal de revisión sistematica.